{{دیوان
| عنوان = تضمین بر شعرِ صائبؔ
| ماخذ = ویکی ماخذ
| ماخذ_لنک = https://ur.wikisource.org/wiki/%D8%AA%D8%B6%D9%85%DB%8C%D9%86_%D8%A8%D8%B1_%D8%B4%D8%B9%D8%B1%D9%90_%D8%B5%D8%A7%D8%A6%D8%A8
| مقابلہ = اقبال سائبر لائبریری (اقبال اکادمی پاکستان)
| مقابلہ_لنک = https://www.iqbalcyberlibrary.net
}}
<poem>
کہاں اقبالؔ تُو نے آ بنایا آشیاں اپنا
نوا اس باغ میں بُلبل کو ہے سامانِ رُسوائی

شرارے وادیِ ایمن کے تُو بوتا تو ہے لیکن
نہیں ممکن کہ پھوٹے اس زمیں سے تُخمِ سِینائی

کلی زورِ نفَس سے بھی وہاں گُل ہو نہیں سکتی
جہاں ہر شے ہو محرومِ تقاضائے خود افزائی

قیامت ہے کہ فطرت سو گئی اہلِ گلستاں کی
نہ ہے بیدار دل پِیری، نہ ہمّت خواہ برنائی

دلِ آگاہ جب خوابیدہ ہو جاتے ہیں سینوں میں
نوا گر کے لیے زہراب ہوتی ہے شکرخائی

نہیں ضبطِ نوا ممکن تو اُڑ جا اس گُلستاں سے
کہ اس محفل سے خوشتر ہے کسی صحرا کی تنہائی

“ہماں بہتر کہ لیلیٰ در بیاباں جلوہ گر باشد
ندارد تنگناے شہر تابِ حُسنِ صحرائی”
</poem>
